\scanpage{165}
Le $\widetilde{\pi}_1(S,\xi)$, $S$ \unsure{préschéma|préschém.} connexe, $\xi$ pt géom.

Pour tout groupe $\pi$, soit
\[
\unsure{\Pi|\pi}^{1}(S,\unsure{\xi|\vert};\pi)
\]
l'ens des classes, à isom près, de \deleted{\unsure{rev.|esp.}} \inserted{fibrés} principaux
homogènes sur $S$, de groupe structural $\pi$,
ponctués au dessus de $\xi$.

(i) $\pi\leadsto\Pi^1(S,\xi;\pi)$ est exact à gauche.

En effet \quad a) Commute aux produits finis
c'est trivial

b) Transforme monos \inserted{$\pi\mathrel{\unsure{-|\to}}\pi'$} en \unsure{mono.|mon.}

En effet, cela provient que \deleted{\unsure{pour tt}} \inserted{si}
\deleted{isom} $P'_1\mathrel{\overset{u'}{\unsure{-|\to}}}P'_2$ \deleted{\unsure{des|de}} \inserted{deux} $\pi'$ fibrés \deleted{\unsure{est}} isom,
il existe un et un seul tel isom
qui applique $\xi'_1$ en $\xi'_2$. Si alors
$P'_1$ et $P'_2$ sont \uline{associés} à $P_1,P_2$ d'où
\[
\left\{\begin{aligned}P_i&\subset P'_i\\ \xi_i&=\xi'_i\end{aligned}\right.,
\]
\deleted{\unsure{donc}} je dis que $\unsure{u|u'}$ \deleted{applique}
\deleted{\unsure{$P_i$}} induit un isom $u:P_1\mathrel{-}P_2$. En
effet, \deleted{$P_2$ est ouvert et fermé dans $P'_2$ et \unsure{contient} $\xi'_2$,}
$\unsure{u^{-1}|{u'}^{-1}}(P_2)$ est \uline{ouvert} et \uline{fermé} dans
$P'_1$, et contient $\xi$, enfin il est \unsure{inv. pour|invariant pour|sous-groupe}
% archive sheet number: 164
